\section{本章小结}
\label{sec:ipspace-summary}

章首的两个问题都靠内积回答。内积的公理与维数无关，复情形约定对第二个变量线性、对第一个变量共轭线性；Cauchy--Schwarz不等式使$\sqrt{\langle f,f\rangle}$成为范数，并使内积连续。内积范数满足平行四边形恒等式，一致范数与$L^1$范数不满足，所以它们不来自任何内积。

Hilbert空间是在诱导范数下完备的内积空间，也就是范数来自内积的Banach空间。有限维时完备性自动成立，无限维时不然：有限支撑数列和平方积分下的连续函数都是不完备的内积空间。$L^2(\mu)$把平方可积函数的等价类收进来，乘积积分在Hölder不等式下有定义，诱导的正是$L^2$范数，由第~\ref{chap:normed-spaces}章的完备性直接成为Hilbert空间；计数测度下得到$\ell^2$。

在Hilbert空间的闭子空间上，最佳逼近存在且唯一：平行四边形恒等式使极小化序列成为柯西列，完备性给出极限，闭性保证极限留在子空间中。最近点由误差与子空间正交刻画，定义为正交投影，给出唯一的正交分解与勾股关系；有限支撑数列在$\ell^2$中的例子说明闭性不能去掉。一般闭凸集上最近点仍存在唯一，但正交条件要换成变分不等式，归卷五处理。

完备正交系只要求有限组合在范数下稠密，与只允许有限组合的代数基不同；它可以用“与所有方向正交的只有零”来检验，可分Hilbert空间的正交基至多可数，并可由Gram--Schmidt正交化得到。用内积系数展开时，有限截断就是正交投影，Bessel不等式控制系数；完备正交系给出依范数收敛的展开、Parseval等式和尾和形式的截断误差，并使空间与$\ell^2$同构；不完备时只恢复闭张成中的部分。范数收敛不等于点态收敛，Legendre部分和在$L^2$中收敛到$\operatorname{sgn}$，却不一致收敛。

这些结构为后面的章节准备了依据。第~\ref{chap:linear-functionals}章用正交投影证明连续线性泛函都能表示为内积，从而用内积表示观测；第~\ref{chap:fourier-transform}章和第~\ref{chap:wavelet-transform}章分别证明三角函数系与小波系的完备性；截断误差的衰减速度与函数光滑性的关系，由第~\ref{chap:function-approximation}章研究。
